\section{Revolución del aprendizaje por refuerzo}

El aprendizaje por refuerzo permite mejorar las respuestas de un modelo de lenguaje a partir de la evaluación de sus propios intentos. Cuando una solución puede comprobarse, como ocurre en determinados problemas matemáticos o tareas de programación, esa verificación proporciona una señal de recompensa para orientar el aprendizaje. El aprendizaje por refuerzo con recompensas verificables (RLVR) convierte así la capacidad de comprobar resultados en una herramienta de entrenamiento~\cite{shao2024deepseekmath,deepseek2025r1}. Un ejemplo de este avance es DeepSeek-R1-Zero, cuyo entrenamiento mediante Group Relative Policy Optimization (GRPO) con recompensas basadas en reglas elevó el desempeño en problemas matemáticos de AIME 2024 del 15,6\,\% al 71,0\,\% en la métrica pass@1, sin una etapa inicial de ajuste fino supervisado~\cite{deepseek2025r1}.

Sin embargo, trasladar estos avances a dominios abiertos plantea una dificultad: ¿cómo orientar el aprendizaje cuando la calidad de una respuesta no puede comprobarse por completo mediante reglas deterministas? Esta pregunta motiva el presente trabajo, que propone una metodología de adaptación basada en rúbricas extraídas automáticamente de respuestas de referencia y en ayudas que se retiran progresivamente durante el entrenamiento.
